\section{Normal rings}
\label{section-normal-rings}

\noindent
We first introduce the notion of a normal domain, and then we
introduce the (very general) notion of a normal ring.

\begin{definition}
\label{definition-domain-normal}
A domain $R$ is called {\it normal} if it is integrally
closed in its field of fractions.
\end{definition}

\begin{lemma}
\label{lemma-integral-closure-in-normal}
Let $R \to S$ be a ring map.
If $S$ is a normal domain, then the integral closure of $R$
in $S$ is a normal domain.
\end{lemma}

\begin{proof}
Omitted.
\end{proof}

\noindent
The following notion is occasionally useful when
studying normality.

\begin{definition}
\label{definition-almost-integral}
Let $R$ be a domain.
\begin{enumerate}
\item An element $g$ of the fraction
field of $R$ is called {\it almost integral over $R$}
if there exists an element $r \in R$, $r\not = 0$
such that $rg^n \in R$ for all $n \geq 0$.
\item The domain $R$ is called {\it completely normal} if every
almost integral element of the fraction field of $R$ is
contained in $R$.
\end{enumerate}
\end{definition}

\noindent
The following lemma shows that a Noetherian domain is normal
if and only if it is completely normal.

\begin{lemma}
\label{lemma-almost-integral}
Let $R$ be a domain with fraction field $K$.
If $u, v \in K$ are almost integral over $R$, then so are
$u + v$ and $uv$. Any element $g \in K$ which is integral over $R$
is almost integral over $R$. If $R$ is Noetherian
then the converse holds as well.
\end{lemma}

\begin{proof}
If $ru^n \in R$ for all $n \geq 0$ and
$v^nr' \in R$ for all $n \geq 0$, then
$(uv)^nrr'$ and $(u + v)^nrr'$ are in $R$ for
all $n \geq 0$. Hence the first assertion.
Suppose $g \in K$ is integral over $R$.
In this case there exists an $d > 0$ such that
the ring $R[g]$ is generated by $1, g, \ldots, g^d$ as an $R$-module.
Let $r \in R$ be a common denominator of the elements
$1, g, \ldots, g^d \in K$. It follows that $rR[g] \subset R$,
and hence $g$ is almost integral over $R$.

\medskip\noindent
Suppose $R$ is Noetherian and $g \in K$ is almost integral over $R$.
Let $r \in R$, $r\not = 0$ be as in the definition.
Then $R[g] \subset \frac{1}{r}R$ as an $R$-module.
Since $R$ is Noetherian this implies that $R[g]$ is
finite over $R$. Hence $g$ is integral over $R$, see
Lemma \ref{lemma-finite-is-integral}.
\end{proof}

\begin{lemma}
\label{lemma-localize-normal-domain}
Any localization of a normal domain is normal.
\end{lemma}

\begin{proof}
Let $R$ be a normal domain, and let $S \subset R$ be
a multiplicative subset. Suppose $g$ is an element
of the fraction field of $R$ which is integral over $S^{-1}R$.
Let $P = x^d + \sum_{j < d} a_j x^j$ be a polynomial
with $a_i \in S^{-1}R$ such that $P(g) = 0$.
Choose $s \in S$ such that $sa_i \in R$ for all $i$.
Then $sg$ satisfies the monic polynomial
$x^d + \sum_{j < d} s^{d-j}a_j x^j$ which has coefficients
$s^{d-j}a_j$ in $R$. Hence $sg \in R$ because $R$ is normal.
Hence $g \in S^{-1}R$.
\end{proof}

\begin{lemma}
\label{lemma-PID-normal}
A principal ideal domain is normal.
\end{lemma}

\begin{proof}
Let $R$ be a principal ideal domain.
Let $g = a/b$ be an element of the fraction field
of $R$ integral over $R$. Because $R$ is a principal ideal domain
we may divide out a common factor of $a$ and $b$
and assume $(a, b) = R$. In this case, any equation
$(a/b)^n + r_{n-1} (a/b)^{n-1} + \ldots + r_0 = 0$
with $r_i \in R$ would imply $a^n \in (b)$. This
contradicts $(a, b) = R$ unless $b$ is a unit in $R$.
\end{proof}

\begin{lemma}
\label{lemma-prepare-polynomial-ring-normal}
Let $R$ be a domain with fraction field $K$.
Suppose $f = \sum \alpha_i x^i$ is an
element of $K[x]$.
\begin{enumerate}
\item If $f$ is integral over $R[x]$
then all $\alpha_i$ are integral over $R$, and
\item If $f$ is almost integral over $R[x]$
then all $\alpha_i$ are almost integral over $R$.
\end{enumerate}
\end{lemma}

\begin{proof}
We first prove the second statement.
Write $f = \alpha_0 + \alpha_1 x + \ldots + \alpha_r x^r$
with $\alpha_r \not = 0$.
By assumption there exists $h = b_0 + b_1 x + \ldots + b_s x^s \in R[x]$,
$b_s \not = 0$ such that $f^n h \in R[x]$ for all
$n \geq 0$. This implies that $b_s \alpha_r^n \in R$
for all $n \geq 0$. Hence $\alpha_r$ is almost
integral over $R$. Since the set of almost integral
elements form a subring (Lemma \ref{lemma-almost-integral}) we deduce that
$f - \alpha_r x^r = \alpha_0 + \alpha_1 x + \ldots + \alpha_{r - 1} x^{r - 1}$
is almost integral over $R[x]$. By induction on $r$ we win.

\medskip\noindent
In order to prove the first statement we will use absolute Noetherian
reduction. Namely, write $\alpha_i = a_i / b_i$ and
let $P(t) = t^d + \sum_{j < d} f_j t^j$ be a polynomial
with coefficients $f_j \in R[x]$ such that $P(f) = 0$.
Let $f_j = \sum f_{ji}x^i$. Consider the subring
$R_0 \subset R$ generated by the finite list of elements
$a_i, b_i, f_{ji}$ of $R$. It is a domain; let
$K_0$ be its field of fractions. Since $R_0$ is a finite type
$\mathbf{Z}$-algebra it is Noetherian, see
Lemma \ref{lemma-obvious-Noetherian}. It is still
the case that $f \in K_0[x]$ is integral over $R_0[x]$,
because all the identities in $R$
among the elements $a_i, b_i, f_{ji}$ also hold in $R_0$.
By Lemma \ref{lemma-almost-integral} the element
$f$ is almost integral over $R_0[x]$. By the second statement of
the lemma, the elements $\alpha_i$ are almost integral
over $R_0$. And since $R_0$ is Noetherian, they are
integral over $R_0$, see Lemma \ref{lemma-almost-integral}.
Of course, then they are integral over $R$.
\end{proof}

\begin{lemma}
\label{lemma-polynomial-domain-normal}
Let $R$ be a normal domain.
Then $R[x]$ is a normal domain.
\end{lemma}

\begin{proof}
The result is true if $R$ is a field $K$ because
$K[x]$ is a euclidean domain and hence a principal ideal
domain and hence normal by Lemma \ref{lemma-PID-normal}.
Let $g$ be an element of the fraction field of
$R[x]$ which is integral over $R[x]$. Because $g$
is integral over $K[x]$ where $K$ is the fraction
field of $R$ we may write $g = \alpha_d x^d + \alpha_{d-1}x^{d-1} +
\ldots + \alpha_0$ with $\alpha_i \in K$.
By Lemma \ref{lemma-prepare-polynomial-ring-normal}
the elements $\alpha_i$ are integral over $R$ and
hence are in $R$.
\end{proof}

\begin{lemma}
\label{lemma-power-series-over-Noetherian-normal-domain}
Let $R$ be a Noetherian normal domain. Then $R[[x]]$ is
a Noetherian normal domain.
\end{lemma}

\begin{proof}
The power series ring is Noetherian by
Lemma \ref{lemma-Noetherian-power-series}.
Let $f, g \in R[[x]]$ be nonzero elements such that
$w = f/g$ is integral over $R[[x]]$.
Let $K$ be the fraction field of $R$. Since the ring of Laurent series
$K((x)) = K[[x]][1/x]$ is a field, we can write
$w = a_n x^n + a_{n + 1} x^{n + 1} + \ldots$
for some $n \in \mathbf{Z}$, $a_i \in K$, and $a_n \not = 0$.
By Lemma \ref{lemma-almost-integral} we see there exists a
nonzero element $h = b_m x^m + b_{m + 1} x^{m + 1} + \ldots$
in $R[[x]]$ with $b_m \not = 0$ such that
$w^e h \in R[[x]]$ for all $e \geq 1$. We conclude that $n \geq 0$ and that
$b_m a_n^e \in R$ for all $e \geq 1$.
Since $R$ is Noetherian this implies that $a_n \in R$ by
the same lemma. Now, if $a_n, a_{n + 1}, \ldots, a_{N - 1} \in R$,
then we can apply the same argument to
$w - a_n x^n - \ldots - a_{N - 1} x^{N - 1} = a_N x^N + \ldots$.
In this way we see that all $a_i \in R$ and the lemma is proved.
\end{proof}

\begin{lemma}
\label{lemma-normality-is-local}
Let $R$ be a domain. The following are equivalent:
\begin{enumerate}
\item The domain $R$ is a normal domain,
\item for every prime $\mathfrak p \subset R$ the local ring
$R_{\mathfrak p}$ is a normal domain, and
\item for every maximal ideal $\mathfrak m$ the ring $R_{\mathfrak m}$
is a normal domain.
\end{enumerate}
\end{lemma}

\begin{proof}
We deduce (1) $\Rightarrow$ (2) from Lemma \ref{lemma-localize-normal-domain}.
The implication (2) $\Rightarrow$ (3) is immediate. The implication
(3) $\Rightarrow$ (1) follows from the fact that for any domain $R$ we have
$$
R = \bigcap\nolimits_{\mathfrak m} R_{\mathfrak m}
$$
inside the fraction field of $R$. Namely, if $g$ is an element of
the right hand side then the ideal $I = \{x \in R \mid xg \in R\}$
is not contained in any maximal ideal $\mathfrak m$, whence $I = R$.
\end{proof}

\noindent
Lemma \ref{lemma-normality-is-local} shows that the following definition
is compatible with Definition \ref{definition-domain-normal}. (It is the
definition from EGA -- see \cite[IV, 5.13.5 and 0, 4.1.4]{EGA}.)

\begin{definition}
\label{definition-ring-normal}
A ring $R$ is called {\it normal} if for every prime
$\mathfrak p \subset R$ the localization $R_{\mathfrak p}$ is
a normal domain (see Definition \ref{definition-domain-normal}).
\end{definition}

\noindent
Note that a normal ring is a reduced ring, as $R$ is a subring of the product
of its localizations at all primes (see for example
Lemma \ref{lemma-characterize-zero-local}).

\begin{lemma}
\label{lemma-normal-ring-integrally-closed}
A normal ring is integrally closed in its total ring of fractions.
\end{lemma}

\begin{proof}
Let $R$ be a normal ring. Let $x \in Q(R)$ be an element of the total ring
of fractions of $R$ integral over $R$. Set $I = \{f \in R, fx \in R\}$. Let
$\mathfrak p \subset R$ be a prime. As $R \to R_{\mathfrak p}$ is
flat we see that $R_{\mathfrak p} \subset Q(R) \otimes_R R_{\mathfrak p}$. As
$R_{\mathfrak p}$ is a normal domain we see that $x \otimes 1$ is an element of
$R_{\mathfrak p}$. Hence we can find $a, f \in R$, $f \not \in \mathfrak p$
such that $x \otimes 1 = a \otimes 1/f$. This means that $fx - a$ maps to
zero in $Q(R) \otimes_R R_{\mathfrak p} = Q(R)_{\mathfrak p}$, which
in turn means that there exists an $f' \in R$, $f' \not \in \mathfrak p$
such that $f'fx = f'a$ in $R$. In other words, $ff' \in I$. Thus $I$
is an ideal which isn't contained in any of the prime ideals of $R$, i.e.,
$I = R$ and $x \in R$.
\end{proof}

\begin{lemma}
\label{lemma-localization-normal-ring}
A localization of a normal ring is a normal ring.
\end{lemma}

\begin{proof}
Omitted.
\end{proof}

\begin{lemma}
\label{lemma-polynomial-ring-normal}
Let $R$ be a normal ring. Then $R[x]$ is a normal ring.
\end{lemma}

\begin{proof}
Let $\mathfrak q$ be a prime of $R[x]$. Set $\mathfrak p = R \cap \mathfrak q$.
Then we see that $R_{\mathfrak p}[x]$ is a normal domain by
Lemma \ref{lemma-polynomial-domain-normal}.
Hence $(R[x])_{\mathfrak q}$ is a normal domain by
Lemma \ref{lemma-localize-normal-domain}.
\end{proof}

\begin{lemma}
\label{lemma-finite-product-normal}
A finite product of normal rings is normal.
\end{lemma}

\begin{proof}
It suffices to show that the product of two normal rings, say $R$ and $S$, is
normal. By Lemma \ref{lemma-disjoint-decomposition} the prime ideals of
$R\times S$ are of the form $\mathfrak{p}\times S$ and $R\times
\mathfrak{q}$, where $\mathfrak{p}$ and $\mathfrak{q}$ are primes of $R$
and $S$ respectively. Localization yields 
$(R\times S)_{\mathfrak{p}\times S}=R_{\mathfrak{p}}$ which is a normal domain
by assumption. Similarly for $S$.
\end{proof}

\begin{lemma}
\label{lemma-characterize-reduced-ring-normal}
Let $R$ be a ring. Assume $R$ is reduced and has finitely many
minimal primes. Then the following are equivalent:
\begin{enumerate}
\item $R$ is a normal ring,
\item $R$ is integrally closed in its total ring of fractions, and
\item $R$ is a finite product of normal domains.
\end{enumerate}
\end{lemma}

\begin{proof}
The implications (1) $\Rightarrow$ (2) and
(3) $\Rightarrow$ (1) hold in general,
see Lemmas \ref{lemma-normal-ring-integrally-closed} and
\ref{lemma-finite-product-normal}.

\medskip\noindent
Let $\mathfrak p_1, \ldots, \mathfrak p_n$ be the minimal primes of $R$.
By Lemmas \ref{lemma-reduced-ring-sub-product-fields} and
\ref{lemma-total-ring-fractions-no-embedded-points} we have
$Q(R) = R_{\mathfrak p_1} \times \ldots \times R_{\mathfrak p_n}$, and
by Lemma \ref{lemma-minimal-prime-reduced-ring} each factor is a field.
Denote $e_i = (0, \ldots, 0, 1, 0, \ldots, 0)$ the $i$th idempotent
of $Q(R)$.

\medskip\noindent
If $R$ is integrally closed in $Q(R)$, then it contains in particular
the idempotents $e_i$, and we see that $R$ is a product of $n$
domains (see Sections \ref{section-connected-components} and
\ref{section-tilde-module-sheaf}). Each factor is of the form
$R/\mathfrak p_i$ with field of fractions $R_{\mathfrak p_i}$. 
By Lemma \ref{lemma-finite-product-integral-closure} each map
$R/\mathfrak p_i \to R_{\mathfrak p_i}$ is integrally closed. 
Hence $R$ is a finite product of normal domains.
\end{proof}

\begin{lemma}
\label{lemma-colimit-normal-ring}
Let $(R_i, \varphi_{ii'})$ be a directed system
(Categories, Definition \ref{definition-directed-system})
of rings. If each $R_i$ is a normal ring so is
$R = \colim_i R_i$.
\end{lemma}

\begin{proof}
Let $\mathfrak p \subset R$ be a prime ideal.
Set $\mathfrak p_i = R_i \cap \mathfrak p$ (usual abuse of notation).
Then we see that
$R_{\mathfrak p} = \colim_i (R_i)_{\mathfrak p_i}$.
Since each $(R_i)_{\mathfrak p_i}$ is a normal domain we
reduce to proving the statement of the lemma for normal
domains. If $a, b \in R$ and $a/b$ satisfies a monic polynomial
$P(T) \in R[T]$, then we can find a (sufficiently large) $i \in I$
such that $a, b$ come from objects $a_i, b_i$ over $R_i$, $P$ comes from a
monic polynomial $P_i\in R_i[T]$ and $P_i(a_i/b_i)=0$. Since $R_i$ is normal we
see $a_i/b_i \in R_i$ and hence also $a/b \in R$.
\end{proof}







